\documentclass[12pt]{article}
\usepackage[utf8]{inputenc}
\usepackage[russian]{babel}
\usepackage{amsmath}
\usepackage{amssymb}
\usepackage{geometry}
\geometry{a4paper, margin=2cm}

\begin{document}

\section*{Задача 1.6}

\textbf{Условие:} Написать программу разворота бинарной
строки в алфавите $\{\#,0,1\}$. Провести анализ сложности
в среднем и в худшем случаях.

\section*{Программа}

\begin{verbatim}
alphabet = [#, 0, 1]

tapes = [ right ]

START:
    # -> #R CheckEmpty
CheckEmpty:
    # -> #H STOP
    0 -> #R Right_0
    1 -> #R Right_1
Right_0:
    # -> 0L Left
    0 -> 0R Right_0
    1 -> 1R Right_0
Right_1:
    # -> 1L Left
    0 -> 0R Right_1
    1 -> 1R Right_1
Left:
    # -> #R RMove
    0 -> #L Left_0
    1 -> #L Left_1
Left_0:
    # -> 0R Right
    0 -> 0L Left_0
    1 -> 1L Left_0
Left_1:
    # -> 1R Right
    0 -> 0L Left_1
    1 -> 1L Left_1
Right:
    # -> #R RMove
    0 -> #R Right_0
    1 -> #R Right_1
RMove:
    # -> #L Shift_#
    0 -> 0R RMove
    1 -> 1R RMove
Shift_#:
    # -> #H STOP
    0 -> #L Shift_0
    1 -> #L Shift_1
Shift_1:
    # -> 1H STOP
    0 -> 1L Shift_0
    1 -> 1L Shift_1
Shift_0:
    # -> 0H STOP
    0 -> 0L Shift_0
    1 -> 0L Shift_1
\end{verbatim}

\subsection*{Сложность}

Пусть $n = |w|$, без \#.

\paragraph{$n=0$.}
\[
T(\varnothing)=2
\]

Далее пусть $n\geq 1$.

В состояниях $Right_0$ и $Right_1$ длины проходов вправо -
$n,n-2,n-4,\ldots$:
\[
R(n)=n+(n-2)+(n-4)+\cdots
=\left\lfloor\frac{(n+1)^2}{4}\right\rfloor
\]

В состояниях $Left_0$ и $Left_1$ длины проходов влево -
$n-1,n-3,n-5,\ldots$:
\[
L(n)=(n-1)+(n-3)+(n-5)+\cdots
=\left\lfloor\frac{n^2}{4}\right\rfloor
\]

\begin{enumerate}
\item В начале два шага: $START\to CheckEmpty\to Right_i$.

\item $Left$ вызывается после каждого прохода вправо
в $Right_0$ и $Right_1$. Их $\left\lceil n/2\right\rceil$.

\item $Right$ вызывается после завершения прохода влево
в $Left_0$ или $Left_1$. Их $\left\lfloor n/2\right\rfloor$.

\item В $RMove$ надо пройти
$\left\lfloor(n+1)/2\right\rfloor$ символов,
а затем выполняет ещё один переход на $\#$ в $Shift_\#$. 
Поэтому $\left\lfloor\frac{n+1}{2}\right\rfloor+1$.

\item $Shift_\#$ - один шаг.

\item После этого $Shift_0$ и $Shift_1$ сдвигают
влево по и завершают работу на левом \#.
Это $\left\lceil n/2\right\rceil$ шагов.
\end{enumerate}

Значит, вообще шагов:
\[
T(w)=2+R(n)+L(n)
+\left\lceil\frac n2\right\rceil
+\left\lfloor\frac n2\right\rfloor
+\left\lfloor\frac{n+1}{2}\right\rfloor+1+1
+\left\lceil\frac n2\right\rceil
\]

\[
\left\lfloor\frac{n^2}{4}\right\rfloor
+\left\lfloor\frac{(n+1)^2}{4}\right\rfloor
=\frac{n(n+1)}2,
\left\lfloor\frac n2\right\rfloor
+\left\lceil\frac n2\right\rceil=n
\]

Значит, сложность:
\[
T(w)=\frac{n(n+1)}2+n+2\left\lceil\frac n2\right\rceil+4
\]

\subsection*{Сложность в худшем}

Число переходов не зависит от значений битов (это видно по программе).
Значит, все слова обрабатываются за одинаковое время.
Ну и можно раскрыть округление вверх:

\[
T_{\max}(n)=
\begin{cases}
2, & n=0\\
\dfrac{n^2+5n}{2}+4, & n\geq 2,\ n\;\vdots\; 2\\
\dfrac{n^2+5n}{2}+5, & n\not\vdots\; 2
\end{cases}
\]

\subsection*{Сложность в среднем}

Для всех слов время работы одинаково $T(w)$, значит
\[
T_{avg}(n)
=\frac{1}{2^n}\sum_{w\in\{0,1\}^n}T(w)
=\frac{2^nT_{\max}(n)}{2^n}
=T_{\max}(n)
\]

Значит, сложность в среднем:
\[
T_{avg}(n)=
\begin{cases}
2, & n=0\\
\dfrac{n^2+5n}{2}+4, & n\geq 2,\ n\;\vdots\; 2\\
\dfrac{n^2+5n}{2}+5, & n\not\vdots\; 2
\end{cases}
\]

\end{document}